\section[Notebook 1971-4: a monoid and its Grothendieck group]{Notebook 1971-4: the group completion of a commutative monoid is a flat functor\protect\verifie{Notebook1971_4}}\label{sec:quillen-1971-4}

\emph{Context.} \carnet{1971}{4} (20 March -- 11 April 1971) studies torsors under a small category $\mathcal{C}$ over a space or a topos, with a view to Mather's theorem. On PDF page~24 (its \q{8}, in the second pagination), Quillen says that a functor $u : \mathcal{C}_1 \to \mathcal{C}_2$ is \emph{strongly cofinal} if $(u_!, u^*)$ is a morphism of presheaf topoi, which amounts to requiring that, for every object $Y$ of $\mathcal{C}_2$, the functor $X \mapsto \operatorname{Hom}(Y, uX)$ be a pro-object of $\mathcal{C}_1$, in other words that its category of elements be cofiltered. He notes that such a $u$ induces isomorphisms on the $R^q\varprojlim$. On page~25 (its \q{9}), he tests the definition on the universal map from a commutative monoid $M$ to its group completion $G$, viewed as a functor between one-object categories: the condition becomes that the category whose objects are the elements of $G$, and whose arrows are given by translations by $M$, is cofiltered, and he checks the two axioms in a few lines.

\emph{Intuition.} An arrow $g \to g'$ is an $m \in M$ such that $mg = g'$. Two elements of $G$ can be written with a common denominator, $g_i = m_i\,b^{-1}$, and are thus translates of $b^{-1}$: this is the axiom on products. Two parallel arrows $m_1, m_2 : g \to g'$ have the same image in $G$; by the construction of the Grothendieck group, $c\,m_1 = c\,m_2$ for some $c \in M$, and precomposing with the arrow $c : c^{-1}g \to g$ equalizes them: this is the axiom on equalizers.

\begin{encadre}
\begin{theoreme}[Quillen, notebook 1971-4, p.~25]\label{thm:quillen-1971-4}
Let $M$ be a commutative monoid, $G(M)$ its Grothendieck group and $u : M \to G(M)$ the canonical map. The induced functor between one-object categories is representably flat: for the unique object $X$, the category $X/u$ is cofiltered. In particular, it is final.
\end{theoreme}
\end{encadre}

\begin{proof}
The objects of $X/u$ are the elements of $G(M)$, and an arrow $g \to g'$ is an $m \in M$ such that $u(m)\,g = g'$. The category is nonempty (it contains $1$).

Every $x \in G(M)$ satisfies $u(b)\,x = u(a)$ for some $a, b \in M$. For $g_1, g_2$, choose $u(b_i)\,g_i = u(a_i)$ and put $g = u(b_1 b_2)^{-1}$; then $u(b_2 a_1)\,g = g_1$ and $u(b_1 a_2)\,g = g_2$.

If $m_1, m_2 : g \to g'$, then $u(m_1)\,g = u(m_2)\,g$, so $u(m_1) = u(m_2)$, so $c\,m_1 = c\,m_2$ for some $c \in M$, since $G(M)$ is the localization of $M$ at all of $M$. The arrow $c : u(c)^{-1} g \to g$ equalizes $m_1$ and $m_2$.

A representably flat functor is final (mathlib, \texttt{final\_of\_representablyFlat}).
\end{proof}

\begin{sloppypar}
\noindent\emph{Verification in Lean.} In \href{\leanblobmain/Notebook1971_4.lean}{\texttt{Notebook1971\_4.\allowbreak{}lean}}, the functor is \texttt{toGroupFunctor}, induced by \texttt{Algebra.GrothendieckGroup.of}; the theorem is \texttt{representablyFlat\_toGroupFunctor}, in the sense of mathlib's \texttt{CategoryTheory.RepresentablyFlat} (all the categories \texttt{StructuredArrow X F} cofiltered), and \texttt{final\_toGroupFunctor}. Quillen's two axioms are isolated as statements about the Grothendieck group: \texttt{exists\_common} (products) and \texttt{exists\_mul\_eq\_mul} (equalizers), with \texttt{exists\_mul\_eq\_of}. The monoid is written multiplicatively, as on the page.
\end{sloppypar}

Not formalized: the equivalence of the definitions on page~24 (morphism of topoi, pro-objects, cofiltered categories of elements), the remark on the $R^q\varprojlim$, and the statements on torsors, which mathlib does not have over a small category.

\litterature{folklore of group completion, found in the literature: flat functors and filtered categories of elements, SGA~4, exposés~I and~IV; Mac~Lane--Moerdijk, \emph{Sheaves in Geometry and Logic} (1992), VII.6. mathlib has \texttt{RepresentablyFlat} and the Grothendieck group, but not this example. It is proposed for mathlib, and the edition does not present it as new.}
